\documentclass[a4print,italian,lof,lot,twoside]{univpmthesis}
\errorcontextlines=9

\RequirePackage[utf8]{inputenc}
\RequirePackage[T1]{fontenc}

\usepackage{lmodern}



%%%%%%%%%%%%%%%%%%%%%%%%%%%%%%%%%%%%%%%%%%%%%%%%%%%%%%%%%%%
% Metadata
%%%%%%%%%%%%%%%%%%%%%%%%%%%%%%%%%%%%%%%%%%%%%%%%%%%%%%%%%%%
\thfaculty{Facolt\`{a} di Ingegneria}
\thcourse{Corso di Laurea in Ingegneria Informatica e dell'Automazione}
\thtitle{VisualSync: Multi-Camera Synchronization}
\thauthor{Davide Acciarri\\ Giacomo Marcucci }
\thadvisor{Rocco Pietrini}

\ayear{2025-2026}

%\thesisdedication{A tutti gli Studenti dell'UNIVPM\\ per i loro primi passi in \LaTeX}
\thlocation{Ancona}
\thtime{Luglio 2026} % data scrittura
%%%%%%%%%%%%%%%%%%%%%%%%%%%%%%%%%%%%%%%%%%%%%%%%%%%%%%%%%%%

\usepackage{amsmath,amsfonts,amssymb,bm} %<- scrittura equazioni
%\usepackage[italian,english]{babel}
\usepackage{graphicx} 	%<- immagini
\usepackage{epstopdf} 	%<- conversione immagini .eps
\usepackage{subfig} 	%<- sottofigure
\usepackage{footnote}	%<- note a piè pagina
\usepackage{tabularx,booktabs,multicol,multirow} %<- tabelle
\usepackage{mathrsfs}
\usepackage{caption}	%<- didascalie
\usepackage{microtype}
\usepackage{gensymb,siunitx}	%<- simboli ed unità di misura
\usepackage{hyperref}
\usepackage{booktabs}
\usepackage{array}
\hypersetup{hidelinks}

\usepackage{glossaries}
\usepackage{verbatim}
\usepackage{listings}
\usepackage{xcolor}
\usepackage{tikz}
\usetikzlibrary{shadows.blur}

\lstdefinestyle{bashstyle}{
    language=bash,
    backgroundcolor=\color{gray!10},
    basicstyle=\ttfamily\small,
    keywordstyle=\color{blue},
    commentstyle=\color{gray},
    stringstyle=\color{orange},
    numbers=left,
    numberstyle=\tiny\color{gray},
    frame=single,
    breaklines=true,
    showstringspaces=false,
    tabsize=2,
    mathescape=false
}

\lstdefinestyle{pythonstyle}{
    language=Python,
    backgroundcolor=\color{gray!10},
    basicstyle=\ttfamily\scriptsize,
    keywordstyle=\color{blue},
    commentstyle=\color{gray},
    stringstyle=\color{orange},
    numbers=left,
    numberstyle=\tiny\color{gray},
    frame=single,
    breaklines=true,
    breakatwhitespace=false,
    showstringspaces=false,
    tabsize=2,
    mathescape=false
}
\usepackage[a-1b]{pdfx} %<- generazione pdf-A

\newcommand{\eps}{\epsilon}
\newcommand{\veps}{\varepsilon}
\newcommand{\Qed}{\begin{flushright}\qed\end{flushright}}

\newcommand{\parinn}{\setlength{\parindent}{1cm}}
\newcommand{\parinf}{\setlength{\parindent}{0cm}}

% \newcommand{\norm}{\|\cdot\|}
\newcommand{\inorm}{\norm_{\infty}}
\newcommand{\opensets}{\{V_{\alpha}\}_{\alpha\in I}}
\newcommand{\oset}{V_{\alpha}}
\newcommand{\opset}[1]{V_{\alpha_{#1}}}
\newcommand{\lub}{\text{lub}}
\newcommand{\del}[2]{\frac{\partial #1}{\partial #2}}
\newcommand{\Del}[3]{\frac{\partial^{#1} #2}{\partial^{#1} #3}}
\newcommand{\deld}[2]{\dfrac{\partial #1}{\partial #2}}
\newcommand{\Deld}[3]{\dfrac{\partial^{#1} #2}{\partial^{#1} #3}}
\newcommand{\der}[2]{\frac{\mathrm{d} #1}{\mathrm{d} #2}}
% \newcommand{\ddd}[3]{\frac{\mathrm{d}^{#3} #1}{\mathrm{d}^{#3} #2}}
\newcommand{\lm}{\lambda}
\newcommand{\uin}{\mathbin{\rotatebox[origin=c]{90}{$\in$}}}
\newcommand{\usubset}{\mathbin{\rotatebox[origin=c]{90}{$\subset$}}}
\newcommand{\lt}{\left}
\newcommand{\rt}{\right}
\newcommand{\bs}[1]{\boldsymbol{#1}}
\newcommand{\exs}{\exists}
\newcommand{\st}{\strut}
\newcommand{\dps}[1]{\displaystyle{#1}}
\newcommand{\id}{\text{id}}

\newcommand{\sess}{\underline{se e solo se} }
\newcommand{\non}{\underline{non} }
\newcommand{\Non}{\underline{Non} }

\newcommand{\ul}[1]{\underline{#1}}
\newcommand{\U}[1]{\underline{#1}}
\newcommand{\bl}[1]{\textbf{#1}}
\newcommand{\B}[1]{\textbf{#1}}
\newcommand{\itl}[1]{\textit{#1}}
\newcommand{\I}[1]{\textit{#1}}

\newcommand{\affianca}[2]{
    \begin{minipage}{0.45\textwidth}
        #1
    \end{minipage}
    \hfill
    \begin{minipage}{0.45\textwidth}
        #2
    \end{minipage}
}

\newcommand{\myaffianca}[4]{
    \begin{minipage}{#1\textwidth}
        #2
    \end{minipage}
    \hfill
    \begin{minipage}{#3\textwidth}
        #4
    \end{minipage}
}

\newcommand{\triaffianca}[3]{
    \begin{minipage}{0.3\textwidth}
        #1
    \end{minipage}
    \hfill
    \begin{minipage}{0.3\textwidth}
        #2
    \end{minipage}
    \hfill
    \begin{minipage}{0.3\textwidth}
        #3
    \end{minipage}
}

\newcommand{\image}[2]{
    \smallskip
    \centering
    \includegraphics[width=#1\textwidth]{#2}
    \smallskip }

\newcommand{\tabimage}[2]{
    \includegraphics[width=#1\textwidth]{#2}
}

\newcommand{\e}{È }


\newcommand{\p}[1]{
    
    \noindent 
    #1}




\input{tools/letterfonts}

%%%%%%%%%%%%%%%%%%%%%%%%%%%%%%%%%%%%%%%%%%%%%%%%%%%%%%%%%%%
\newcommand{\eqname}{Eq.}
%%%%%%%%%%%%%%%%%%%%%%%%%%%%%%%%%%%%%%%%%%%%%%%%%%%%%%%%%%%

\begin{document}

%%%%%%%%%%%%%%%%%%%%%%%%%%%%%%%%%%%%%%%%%%%%%%%%%%%%%%%%%%%
% Front matter contents
%%%%%%%%%%%%%%%%%%%%%%%%%%%%%%%%%%%%%%%%%%%%%%%%%%%%%%%%%%%
\frontmatter

\makesofttitle

\begin{comment}
    \begin{thesisacknowledge}[italian]
    \end{thesisacknowledge}


\begin{thesisabstract}[italian]
	L'obiettivo di questo studio è la ricerca di metodologie e algoritmi di machine learning che impiegano reti neurali applicabili a dati generati da materassi intelligenti dotati di sensori di pressione piezoresistivi.
    Questo lavoro è parte di un progetto che mira allo sviluppo di una soluzione per la gestione degli ospiti delle Residenze Speciali per Anziani (RSA) che utilizzi dispositivi e sensori all'avanguardia per fornire un valido supporto allo staff nel garantire un elevato standard qualitativo di assistenza, massimizzandone l'efficacia e l’efficienza (sia clinica che operativa) salvaguardando la conformità alle linee guida fornite dalle normative.

    \vspace{8mm}
    
    In questo elaborato ci si è posti l'obiettivo di individuare gli algoritmi e le metodologie presenti in letteratura dedicati al monitoraggio dei pazienti, al fine di migliorare la futura definizione dei requisiti del sistema e la successiva attività di sviluppo.

    \vspace{8mm}

    In conclusione, saranno discusse le prospettive di sviluppo futuro e le possibili evoluzioni del sistema.

    
\end{thesisabstract}
\end{comment}

\thesistoc

%%%%%%%%%%%%%%%%%%%%%%%%%%%%%%%%%%%%%%%%%%%%%%%%%%%%%%%%%%%
% Main matter contents
%%%%%%%%%%%%%%%%%%%%%%%%%%%%%%%%%%%%%%%%%%%%%%%%%%%%%%%%%%%
\mainmatter

%\input{chapters/chapter0/guida}

%\input{chapters/chapter0/chapter0}
\graphicspath{{chapters/chapter1/img}}
\chapter{Introduzione} \label{introduzione}
Con la diffusione di smartphone e dispositivi portatili, è comune che uno stesso evento dinamico, come una partita di calcio, un concerto o una festa in famiglia, venga ripreso da più angolazioni da diverse persone. Tuttavia, allineare questi flussi video per scopi di ricostruzione 4D o editing professionale rimane una sfida complessa a causa della natura frammentaria e casuale delle riprese. 
Per la sincronizzazione temporale di video multi-camera acquisiti in modo indipendente e non coordinato, è stato progettato un framework di ottimizzazione, VisualSync.
I metodi di sincronizzazione esistenti si basano su ambienti controllati, annotazioni manuali, pattern specifici come human pose, segnali audio, o costosi setup hardware come dispositivi time-coded, nessuno dei quali è disponibile in video registrati casualmente. 
Un altro aspetto da dover affrontare e che i video sono spesso "unposed", ovvero non si hanno informazioni a priori riguardanti la posizione spaziale, l'orientamento o i parametri intrinseci delle telecamere al momento della ripresa.
Inoltre, costruire un sistema pratico e robusto che funzioni su video reali è impegnativo, in quanto gli oggetti dinamici sono a priori sconosciuti e generalmente piccoli e sfocati, e non tutte le viste possono avere sovrapposizione.
Il problema parte da un insieme di N video asincroni che riprendono la stessa scena dinamica da diversi punti di
vista e l'obiettivo è sincronizzarli e recuperare un timestamp allineato globalmente. Formalmente, si cerca di stimare un offset temporale per ogni video da applicare al suo tempo di clock originale fuori sincronizzazione. Dopo la sincronizzazione, i frame che condividono lo stesso tempo di clock corrisponderanno allo stesso momento in tutti i video. È possibile vedere  la sincronizzazione globale come un problema di minimizzazione dell'energia ovvero, miriamo a trovare gli offset che allineino meglio tutte le coppie di video minimizzando il loro errore di sincronizzazione pairwise. La minimizzazione dell'energia pone tre sfide: il problema di ottimizzazione è
altamente non convesso, la formulazione è a tempo continuo, ma le osservazioni arrivano a tempi di frame discreti e che negli scenari reali, è difficile associare traiettorie dense tra fotocamere con punti di vista molto diversi e stimare pose accurate per fotocamere in movimento.
Per scrivere questa relazione abbiamo usato come riferimento le seguenti fonti: \cite{sam2}, \cite{visualsync}, \cite{deva}, \cite{uni4d}, \cite{mast3r}, \cite{vggt}, \cite{cotracker3}, \cite{cotracker}
\graphicspath{{chapters/chapter2/img}}
\chapter{Capitolo 2} \label{chapter2}
Lo stato dell'arte relativo alla sincronizzazione e alla ricostruzione di scene dinamiche multi-camera ha subito una trasformazione radicale grazie all'avvento dei modelli di fondazione visiva e di nuove tecniche di ottimizzazione geometrica. La sincronizzazione video è stata esplorata da più prospettive. I metodi basati sulla geometria, come quelli di Albl et al. e Li et al., stimano offset temporali usando la geometria epipolare, ma tipicamente assumono scene statiche o punti di vista fissi. Gli approcci human-centric usano la pose umana come segnale di sincronizzazione, ma sono limitati dall'accuratezza della stima della pose, dal numero di persone nella scena, e faticano in scenari senza attività umana prominente. Gli approcci basati sull'audio usano indizi audio per la sincronizzazione, ma funzionano solo in ambienti silenziosi e non sono
generalmente applicabili a contesti reali in cui l'audio è rumoroso o non disponibile. I metodi basati
sull'apprendimento come Sync-NeRF, ottimizzano congiuntamente pose e offset temporali, ma
spesso sono vincolati ad ambienti o tipi di oggetti specifici. Per la Multi-View Structure-from-Motion ci sono tecniche Structure-from-Motion (SfM), come COLMAP, che hanno significativamente avanzato le pipeline di ricostruzione 3D producendo pose accurate delle fotocamere da immagini e video multi-vista. Modelli più recenti come HLOC, Hierarchical-Localization, e VGGT, Visual Geometry Grounded Transformer, si basano su questi progressi usando feature basate sull'apprendimento e transformer per gestire il matching su larga scala e la stima della pose. Sebbene questi metodi raggiungano forti prestazioni nella stima della geometria delle fotocamere, sono carenti nella sincronizzazione di scene dinamiche, poiché si basano principalmente su indizi visivi statici. Per il tracking e le corrispondenze ci sono modelli come CoTracker tracciano punti densamente nel tempo, offrendo una forte coerenza temporale. Tuttavia, non modellano le corrispondenze spaziali tra punti di vista diversi. D'altro canto, MASt3R si concentra sul matching spaziale e la ricostruzione stereo, fornendo corrispondenze cross-view dense, ma non gestisce la dinamica temporale, specialmente in scene in movimento. La segmentazione si è spostata dai modelli addestrati su classi specifiche come persone o veicoli, a sistemi universali. Il modello SAM 2, Segment Anything Model 2, è l'attuale punto di riferimento ed è un modello fondamentale per la risoluzione del problema della segmentazione visiva guidata da prompt in immagini e video tramite una memoria di streaming che permette di mantenere la coerenza temporale anche in sequenze lunghe e complesse. Framework come DEVA hanno ulteriormente raffinato questo approccio, introducendo un sistema di propagazione bidirezionale che de-normalizza gli errori dei modelli a frame singolo.
\graphicspath{{chapters/chapter3/img}}
\chapter{Strumenti e architetture} \label{strumenti e architetture}
\section{Struttura Dataset}
La struttura del dataset è organizzata in modo gerarchico per gestire le diverse scene di video e i relativi punti di vista. Le diverse scene vengono rappresentate da 18 sottocartelle nominate per ID e all'interno di ogni sottocartella sono presenti tre sottocartelle specifiche per le diverse angolazioni di ripresa: vista dall'alto (TOP), vista in prima persona (FPV) e vista in terza persona (TPV).
Oltre ai video originali, queste cartelle contengono anche i video sincronizzati che forniscono il Ground Truth da utilizzare poi nelle operazioni di confronto.
All'interno del dataset è importante andare a distinguere la gestione delle telecamere in base al loro profilo di movimento per ottimizzare la stima della posa. Le viste TOP e TPV vengono configurate come telecamere statiche. Per idetificare queste due tipologie di visuali statiche, è necessario includere la stringa 'cam' nella nomenclatura delle cartelle durante la fase di preprocessing in cui avviene la frammentazione del video da formato MP4 in immagini in formato JPG.
Al contrario, la vista FPV è classificata come telecamera dinamica poiché, essendo indossata da un soggetto in attività, è soggetta a spostamenti continui e rotazioni repentine. In questo caso non deve esserre inclusa la stringa 'cam' nella nomenclatura delle cartelle durante la fase di preprocessing.
Prima di iniziare l'attività di sincronizzazione dei video, come accennato in precedenza, è stata necessaria una fase di preprocessing. Il framework non elabora direttamente i file video, ma richiede una trasformazione preventiva in sequenze di immagini. Questa trasformazione è stata realizzata attraverso script di preparazione dedicati, in cui i video originali del dataset vengono decodificati e i singoli frame estratti sono riorganizzati in maniera gerarchica. 
\section{Modelli}
\subsection{GPT}
Il modello GPT, Generative Pre-trained Transformer, produce enormi quantità di testo pertinente e complesso generato dalla macchina a partire da una piccola quantità di testo come input. I modelli GPT possono essere identificati come un modello linguistico che imita il testo umano utilizzando tecniche di DL, e agiscono come un modello autoregressive in cui il valore attuale si basa sul valore precedente. La versione del modello GPT integrata all'interno del framework VisualSync è l'ultima, ovvero GPT-4o. GPT-4o è un 'autoregressive Omni model', che accetta come input qualsiasi combinazione di testo, audio, immagini e video e genera qualsiasi combinazione di output di testo, audio e immagini. Il termine 'autoregressive', indica che il modello ha la capacità di prevedere il valore successivo basandosi sui valori precedenti, mentre Omni, da cui deriva la 'o' di GPT-4o, indica che questa capacità di prevedere l'elemento successivo non è limitata al solo testo, ma si estende a qualsiasi tipo di media. È addestrato end-to-end attraverso testo, vision e audio, il che significa che tutti gli input e gli output sono elaborati dalla stessa rete neurale.\cite{GPT-4o} 
L'architettura del modello GPT si basa sul modello transformer. Il modello Transformer utilizza meccanismi di self-attention per elaborare sequenze di input di lunghezza variabile, rendendolo ben adatto per compiti di NLP. GPT semplifica l'architettura sostituendo i blocchi encoder-decoder con blocchi decoder. Il modello GPT prende il modello transformer e lo pre-addestra su grandi quantità di dati di testo. Il processo di pre-addestramento prevede la previsione della parola successiva in una sequenza, date le parole precedenti, un'attività nota come modellazione linguistica. Questo processo di pre-addestramento consente al modello di apprendere rappresentazioni del linguaggio naturale che possono essere sintonizzate per specifici compiti a valle.\cite{GPT}
\subsection{Grounded-SAM 2}
Grounded SAM 2 è un framework che combina Grounding DINO, un modello di object detection open-vocabulary in grado di localizzare oggetti in un'immagine a partire da una descrizione testuale, con SAM 2, Segment Anything Model 2, modello di segmentazione esteso al dominio video grazie a un meccanismo di memoria temporale che ne consente il tracking coerente tra frame. La combinazione dei due moduli permette di ottenere maschere di segmentazione precise per oggetti specificati in linguaggio naturale, senza necessità di fine-tuning su classi predefinite, rendendo il modello particolarmente adatto a scenari multi-camera dove gli oggetti di interesse possono variare da sequenza a sequenza.
\subsubsection{Grounding DINO}
Il modello Grounding DINO è un sistema solido che è stato sviluppato per rilevare oggetti arbitrari specificati da input in linguaggio umano, un compito comunemente denominato 'open-set object detection'. Questo sistema è stato progettato basandosi su due principi: fusione stretta delle modalità basata su DINO, ovvero una fusione stretta tra modalità visiva e testuale, costruita sulla base dell'architettura DINO, e pre-training grounded su larga scala per la generalizzazione di concetti mai visti esplicitamente come classe.
Grounding DINO, dato in input una coppia immagine-testo, restituisce molteplici coppie bounding box-etichetta, dove l'etichetta corrisponde al sintagma nominale estratto dal testo in input. Grounding DINO è un'architettura dual-encoder-single-decoder: contiene un backbone per l'immagine per l'estrazione delle caratteristiche dell'immagine, un backbone per il testo per l'estrazione delle caratteristiche del testo, un potenziatore di caratteristiche per la fusione delle caratteristiche di immagine e testo, un modulo di selezione delle query guidato dal linguaggio per l'inizializzazione delle query e un decodificatore cross-modalità per il perfezionamento del riquadro. Per ogni coppia immagine-testo, vengono estratte caratteristiche standard dell'immagine e caratteristiche standard del testo utilizzando rispettivamente una backcone per le immagini e una per il testo. Le due caratteristiche standard vengono quindi inserite in un modulo di potenziamento delle caratteristiche per la fusione delle caratteristiche tra le diverse modalità. Dopo aver ottenuto le caratteristiche tra testo e immagine tra le diverse modalità, utilizziamo un modulo di selezione delle query guidato dal linguaggio per selezionare le query tra le diverse modalità dalle caratteristiche dell'immagine. Analogamente alle query sugli oggetti, nella maggior parte dei modelli di tipo DETR, Detection Transformers, queste query tra le diverse modalità vengono inserite in un decodificatore per le diverse modalità per analizzare le caratteristiche desiderate dalle due modalità e aggiornarsi automaticamente. Le query di output dell'ultimo livello del decodificatore vengono utilizzate per predire i riquadri degli oggetti ed estrarre le frasi corrispondenti. \cite{GroundingDINO}
\subsubsection{SAM 2}
Il Segment Anything Model 2 è un modello unificato per la segmentazione di video e immagini dove l'immagine è un video a singolo frame. Il modello produce maschere di segmentazione dell'oggetto di interesse, sia su singole immagini sia attraverso i frame video. SAM 2 è dotato di una memoria che immagazzina informazioni sull'oggetto e sulle interazioni precedenti, il che gli consente di generare predizioni di maschere spazio-temporali lungo tutto il video, e anche di correggerle efficacemente sulla base del contesto di memoria immagazzinato relativo all'oggetto, proveniente dai frame osservati in precedenza.
SAM 2 può essere visto come una generalizzazione di SAM, infatti il modello ha un comportamento simile: un decoder di maschere promptable e leggero prende in input un embedding dell'immagine e dei prompt e restituisce una maschera di segmentazione per il frame. L'embedding del frame utilizzato dal decoder di SAM 2 non proviene direttamente da un encoder di immagini, ma è invece condizionato sulle memorie delle predizioni passate e dei frame con prompt. È possibile che i frame con prompt provengano anche 'dal futuro' rispetto al frame corrente. Le memorie dei frame vengono create dal memory encoder sulla base della predizione corrente e collocate in una memory bank per l'uso nei frame successivi. L'operazione di memory attention prende l'embedding per-frame, proveniente dall'image encoder, e lo condiziona sulla memory bank, prima che il mask decoder lo elabori per formare una predizione.
Il modello è composto da diverse componenti:
\begin{itemize}
    \item Image encoder: Per l'elaborazione in tempo reale di video di lunghezza arbitraria, adottiamo un approccio streaming, consumando i frame video man mano che diventano disponibili. L'image encoder viene eseguito una sola volta per l'intera interazione, e il suo ruolo è fornire token non condizionati che rappresentano ciascun frame.
    \item Memory attention: Il ruolo della memory attention è condizionare le feature del frame corrente sulle feature e predizioni dei frame passati, oltre che su eventuali nuovi prompt. Vengono disposti L blocchi transformer, il primo dei quali prende come input l'encoding dell'immagine del frame corrente. Ogni blocco esegue un self-attention, seguito da una cross-attention verso le memorie dei frame e verso i puntatori d'oggetto immagazzinati in una memory bank, ed infine un Multilayer Perceptron, MLP.
    \item Prompt encoder e mask decoder: Il prompt encoder può essere attivato tramite click, box o maschere, per definire l'estensione dell'oggetto in un dato frame. I prompt sparsi sono rappresentati da encoding posizionali sommati a embedding appresi per ciascun tipo di prompt, mentre le maschere vengono incorporate tramite convoluzioni e sommate all'embedding del frame.
    Il design del decoder vede sovrapposti blocchi transformer "two-way" che aggiornano gli embedding dei prompt e del frame. Come in SAM, per prompt ambigui, prediciamo più maschere. Questo design è importante per garantire che il modello produca maschere valide. Nel video, dove l'ambiguità può estendersi attraverso più frame, il modello predice più maschere su ciascun frame. Se nessun prompt successivo risolve l'ambiguità, il modello propaga solo la maschera con l'IoU predetto più alto per il frame corrente.
    A differenza di SAM, dove esiste sempre un oggetto valido da segmentare dato un prompt positivo, è possibile che non esista alcun oggetto valido su alcuni frame e per questa nuova modalità di output, aggiungiamo una head che predice se l'oggetto di interesse è presente nel frame corrente o meno. Un'altra novità sono le connessioni skip dall' image encoder gerarchico per incorporare embedding ad alta risoluzione per il decoding delle maschere.
    \item Memory encoder: Il memory encoder genera una memoria sottocampionando la maschera di output tramite un modulo convoluzionale e sommandola elemento per elemento con l'embedding non condizionato del frame proveniente dall'image encoder, seguito da layer convoluzionali leggeri per fondere l'informazione.
    \item Memory bank: La memory bank conserva informazioni sulle predizioni passate relative all'oggetto target nel video, mantenendo una coda FIFO di memorie fino a N frame recenti, e immagazzina informazioni provenienti dai prompt in una coda FIFO fino a M frame con prompt. Oltre alla memoria spaziale, vengono memorizzati una lista di puntatori d'oggetto come vettori leggeri per l'informazione semantica ad alto livello dell'oggetto da segmentare, basati sui token di output del mask decoder di ciascun frame. La nostra memory attention esegue cross-attention sia verso le feature di memoria spaziale sia verso questi puntatori d'oggetto. Vengono inserite informazioni di posizione temporale nelle memorie degli N frame recenti, permettendo al modello di rappresentare il movimento a breve termine dell'oggetto, ma non in quelle dei frame con prompt, perchè il segnale di addestramento proveniente dai frame con prompt è più scarso ed è più difficile generalizzare al contesto di inferenza, dove i frame con prompt possono provenire da un intervallo temporale molto diverso rispetto a quello del training. \cite{sam2}
\end{itemize} 
\subsection{DEVA}

\subsection{CoTracker3}

\subsection{Mast3R}

\subsection{VGGT}

\section{Le tre fasi di visualsync}
\subsection{Fase 0}
\subsection{Fase 1}
\subsection{Fase 2}
\graphicspath{{chapters/chapter4/img}}
\chapter{Miglioramenti e ottimizzazioni architetturali} \label{miglioramenti}

\graphicspath{{chapters/chapter5/img}}
\chapter{} \label{chapter5}

\graphicspath{{chapters/chapter6/img}}
\chapter{Conclusioni} \label{conclusioni}

\graphicspath{{chapters/chapter7/img}}
\chapter{Possibili sviluppi futuri} \label{sviluppi_futuri}


%%%%%%%%%%%%%%%%%%%%%%%%%%%%%%%%%%%%%%%%%%%%%%%%%%%%%%%%%%%
% Back matter contents
%%%%%%%%%%%%%%%%%%%%%%%%%%%%%%%%%%%%%%%%%%%%%%%%%%%%%%%%%%%
\backmatter

%\bibliographystyle{IEEEbib}
%\bibliography{IEEEabrv,bibliografia}


\bibliographystyle{unsrt}
\bibliography{frontbackmatters/bibliography.bib}

\end{document}
